% !TEX root = ../main.tex
\section{EndorseRank and AWP Side by Side}
\label{ch:appendix-er-awp}

Chapter~\ref{ch:results} reports every score of the thesis together and reads each PageRank against the raw degree it smooths. This appendix collects the direct comparisons between the two single-layer scores, EndorseRank and AWP: their family means and paired differences in the same window, their per-proxy coefficients, their paired differences out of window, and their contrasts on every trader label. None of the comparisons on the matched cohort or on the spender cohorts was fixed in advance, and none of them carries a claim of the thesis; the trader pairs were registered with W1 (Section~\ref{sec:er-awp-neutral-grid}). The two scores walk on different edges and are built for different constructs, so on each check over a full cohort the score ahead is the one whose edges the check counts. Outcomes built from neither edge type are the exception, since neither score was built for them; Section~\ref{sub:neutral-label-results} reports them.

Coverage differs between the two. Of the matched cohort, $962$ contracts lie outside the allowance graph and $3{,}297$ outside the transfer graph, and $3{,}753$ of the $33{,}101$ spenders of W0 had no transfer edge before the freeze, so AWP scores them zero. EndorseRank, by contrast, leaves almost no pair of cohort contracts tied (Section~\ref{sec:rank-divergence}). Out of window the two scores stand differently against the raw degrees they smooth: AWP ranks later new senders above the transfer in-degree in both windows, while EndorseRank ranks later new approval pairs below the in-approve degree in both (Section~\ref{sec:holdout-results}). A lead of EndorseRank over AWP therefore does not show that propagation over approvals adds anything to the count of approving owners.

\dissertationSubheading[sec:er-awp-same-window]{Same window}

Table~\ref{tab:alignment-hybrid} gives the family means of the two scores with their intervals, and the last row of Table~\ref{tab:tie-sensitivity} their agreement. The two patterns cross. AWP tracks the transfer family most closely ($0.447$, interval $[0.440,0.453]$) and the allowance family least ($0.173$, $[0.165,0.180]$), while EndorseRank tracks the allowance family most closely ($0.481$, $[0.476,0.487]$) and stands at $0.122$ ($[0.113,0.129]$) on the transfer family. Table~\ref{tab:tau-diff} tests five paired differences on the same contract resamples (Section~\ref{sec:statistical-tests}), numbered i to v in the order of its rows.

% Generated by experiments/scripts/export_latex.py -- do not edit by hand.
% Source: experiments/data/2-processed-tables-and-evaluations/same-window/eval_summary.json
\begin{table}[htbp]
\centering
\caption[Same-window paired contrasts between EndorseRank and AWP.]{Paired bootstrap tests of differences between Kendall $\tau$ coefficients of EndorseRank and AWP in the same window (matched cohort, $n=27{,}844$ contracts). Family entries use the family-mean $\tau$; the same contract resamples are used for $a$ and $b$, so the interval is on the paired difference. An interval excluding 0 indicates a difference not attributable to sampling variation at the 5\% level. CI = 95\% percentile bootstrap over contracts (400 paired resamples, seed 42).}
\label{tab:tau-diff}
\fitwidth{%
\begin{tabular}{lccccc}
\toprule
Contrast ($a$ minus $b$) & $\tau_a$ & $\tau_b$ & $\Delta\tau$ & 95\% CI of $\Delta\tau$ & $P(\Delta\tau>0)$ \\
\midrule
EndorseRank allowance minus EndorseRank transfer & 0.481 & 0.122 & +0.359 & [0.351, 0.367] & 1.000 \\
AWP transfer minus EndorseRank transfer & 0.447 & 0.122 & +0.325 & [0.315, 0.335] & 1.000 \\
EndorseRank allowance minus AWP allowance & 0.481 & 0.173 & +0.309 & [0.300, 0.317] & 1.000 \\
AWP sybil-stability minus EndorseRank sybil-stability & 0.311 & 0.025 & +0.285 & [0.277, 0.294] & 1.000 \\
EndorseRank in-approve degree minus EndorseRank in-degree & 0.626 & 0.196 & +0.430 & [0.421, 0.440] & 1.000 \\
\bottomrule
\end{tabular}
}
\end{table}


All five intervals lie above zero. EndorseRank's allowance coefficient exceeds its transfer coefficient (contrast~i, $+0.359$, interval $[0.351,0.367]$); AWP's transfer coefficient exceeds EndorseRank's (contrast~ii, $+0.325$, $[0.315,0.335]$); EndorseRank's allowance coefficient exceeds AWP's (contrast~iii, $+0.309$, $[0.300,0.317]$); AWP leads on the Sybil-stability family (contrast~iv, $+0.285$, $[0.277,0.294]$); and EndorseRank follows the in-approve degree more closely than the in-degree (contrast~v, $+0.430$, $[0.421,0.440]$). Contrasts~i and~v compare EndorseRank with itself across proxies. EndorseRank ties only $0.15\%$ of the contract pairs, so its own ties cap none of these coefficients, and the ties of the proxies do not explain gaps of this size: on the in-degree AWP reaches $0.541$ ($[0.534,0.547]$), where EndorseRank reaches $0.196$ ($[0.187,0.204]$). Keeping the contracts outside the allowance graph as isolated nodes changes none of the family means (Table~\ref{tab:tie-sensitivity}).

The opposite signs of contrasts~ii and~iii follow from the constructs. Each score is ahead on the family built from its own edges, and each of these families includes the in-degree that the corresponding PageRank smooths. The agreement is partly built in and checks the intended construct of each score; it is not evidence from outside the score. Every contract of the cohort received approvals from at least three owners, and AWP's allowance coefficients lie above zero (Table~\ref{tab:alignment-allowance}): contracts that many owners approve also tend to be central in the transfer graph. Neither difference says which score is the better measure of reputation.

On the scaling stages of Table~\ref{tab:robustness-sample-size}, whose values are not comparable with those of the full cohort, the order of the two scores within each family is the same at every stage: AWP is ahead on the transfer and Sybil-stability families and EndorseRank on the allowance family.

\dissertationSubheading[sec:proxy-family-detail]{Proxy detail}

Each family mean condenses several per-proxy coefficients. Tables~\ref{tab:alignment-transfer} to~\ref{tab:alignment-sybil-stability} report them one proxy at a time for EndorseRank and AWP, as Spearman $\rho$, Kendall $\tau$ and the 95\% interval of $\tau$. Before correlation, every proxy is oriented so that a higher value means higher expected reputation, and the transfer proxies are the ones used in the evaluation of AWP by its authors \parencite{do2023awp}.

\dissertationSubsubheading[sub:transfer-family]{Transfer family}

Table~\ref{tab:alignment-transfer} gives the coefficients of the two scores with the in-degree and the in-value.

% Generated by experiments/scripts/export_latex.py -- do not edit by hand.
% Source: experiments/data/2-processed-tables-and-evaluations/same-window/eval_summary.json
\begin{table}[htbp]
\centering
\caption[Domain alignment of EndorseRank and AWP with the transfer proxies.]{Domain alignment of EndorseRank and AWP with the transfer proxies (in-degree and in-value of each contract), same window (matched cohort, $n=27{,}844$ contracts). Kendall $\tau$ and Spearman $\rho$ between scores and proxy rankings. CI = 95\% percentile bootstrap over contracts (400 paired resamples, seed 42).}
\label{tab:alignment-transfer}
\fitwidth{%
\begin{tabular}{lccc}
\toprule
Proxy & Spearman $\rho$ & Kendall $\tau$ & 95\% CI of $\tau$ \\
\midrule
\multicolumn{4}{l}{\textit{EndorseRank}} \\
In-degree & 0.283 & 0.196 & [0.187, 0.204] \\
In-value (sum inbound) & 0.081 & 0.048 & [0.039, 0.055] \\
\addlinespace
\multicolumn{4}{l}{\textit{AWP}} \\
In-degree & 0.706 & 0.541 & [0.534, 0.547] \\
In-value (sum inbound) & 0.496 & 0.353 & [0.345, 0.361] \\
\bottomrule
\end{tabular}
}
\end{table}


AWP agrees with the in-degree at $0.541$ (interval $[0.534,0.547]$) and with the in-value at $0.353$ ($[0.345,0.361]$). AWP's value weight saturates near ten base units (Section~\ref{sub:awp-graph}), so the size of a transfer barely enters its score, while the in-value sums raw amounts. EndorseRank's transfer coefficients lie above zero, at $0.196$ ($[0.187,0.204]$) and $0.048$ ($[0.039,0.055]$), and their intervals lie far below AWP's: contracts that many owners approve receive transfers from somewhat more senders than the rest of the cohort. Out of window AWP ranks later new senders above the raw in-degree in both windows (Section~\ref{sec:holdout-results}), so on the transfer side the walk adds to the count it smooths.

\dissertationSubsubheading[sub:allowance-family]{Allowance family}

Table~\ref{tab:alignment-allowance} gives the coefficients of the two scores with the in-approve degree and the in-approve value.

% Generated by experiments/scripts/export_latex.py -- do not edit by hand.
% Source: experiments/data/2-processed-tables-and-evaluations/revised-usd/same-window/eval_summary.json
\begin{table}[htbp]
\centering
\caption[Domain alignment of EndorseRank and AWP with the allowance proxies.]{Domain alignment of EndorseRank and AWP with the allowance proxies (in-approve degree and value of each spender contract), same window (matched cohort, $n=15{,}107$ contracts). Kendall $\tau$ and Spearman $\rho$ between scores and proxy rankings. CI = 95\% percentile bootstrap over contracts (400 paired resamples, seed 42).}
\label{tab:alignment-allowance}
\fitwidth{%
\begin{tabular}{lccc}
\toprule
Proxy & Spearman $\rho$ & Kendall $\tau$ & 95\% CI of $\tau$ \\
\midrule
\multicolumn{4}{l}{\textit{EndorseRank}} \\
In-approve degree & 0.811 & 0.655 & [0.648, 0.662] \\
In-approve value (sum) & 0.681 & 0.515 & [0.506, 0.524] \\
\addlinespace
\multicolumn{4}{l}{\textit{AWP}} \\
In-approve degree & 0.286 & 0.206 & [0.193, 0.217] \\
In-approve value (sum) & 0.260 & 0.180 & [0.169, 0.191] \\
\bottomrule
\end{tabular}
}
\end{table}


EndorseRank agrees with the in-approve degree at $0.626$ (interval $[0.620,0.632]$) and with the in-approve value at $0.336$ ($[0.329,0.344]$). Under the bounded value weight nearly every allowance counts the same, so EndorseRank follows the number of approving owners more than the amounts they approve. For a spender whose owners receive no approvals it equals the closed form of Section~\ref{sub:theory-closed-form}, a function of the owner-normalized in-strength. AWP agrees with the two proxies at $0.227$ ($[0.218,0.235]$) and $0.118$ ($[0.110,0.126]$). The comparison is partly built in, since EndorseRank is computed from the same approvals as the proxies.

\dissertationSubsubheading[sub:sybil-family]{Sybil-adjusted stability family}

Table~\ref{tab:alignment-sybil-stability} gives the coefficients of the two scores with the inbound counterparty ratio, the transfer tenure and the active months.

% Generated by experiments/scripts/export_latex.py -- do not edit by hand.
% Source: experiments/data/2-processed-tables-and-evaluations/revised-usd/same-window/eval_summary.json
\begin{table}[htbp]
\centering
\caption[Domain alignment of EndorseRank and AWP with the Sybil-adjusted stability proxies.]{Domain alignment of EndorseRank and AWP with the Sybil-adjusted stability proxies from transfer activity, same window (matched cohort, $n=15{,}107$ contracts). Kendall $\tau$ and Spearman $\rho$ between scores and proxy rankings. CI = 95\% percentile bootstrap over contracts (400 paired resamples, seed 42).}
\label{tab:alignment-sybil-stability}
\fitwidth{%
\begin{tabular}{lccc}
\toprule
Proxy & Spearman $\rho$ & Kendall $\tau$ & 95\% CI of $\tau$ \\
\midrule
\multicolumn{4}{l}{\textit{EndorseRank}} \\
Inbound counterparty ratio & -0.231 & -0.156 & [-0.166, -0.147] \\
Transfer tenure (days) & 0.124 & 0.080 & [0.069, 0.091] \\
Active months & 0.146 & 0.098 & [0.087, 0.109] \\
\addlinespace
\multicolumn{4}{l}{\textit{AWP}} \\
Inbound counterparty ratio & 0.234 & 0.127 & [0.114, 0.141] \\
Transfer tenure (days) & 0.683 & 0.508 & [0.499, 0.517] \\
Active months & 0.704 & 0.549 & [0.540, 0.558] \\
\bottomrule
\end{tabular}
}
\end{table}


The table shows AWP ahead of EndorseRank on the transfer tenure ($0.452$, interval $[0.445,0.458]$, against $0.083$, $[0.074,0.090]$) and on the active months ($0.499$, $[0.492,0.506]$, against $0.121$, $[0.113,0.129]$), as one would expect from a transfer score that rewards activity spread over the window; the paired difference on the family mean excludes zero (Table~\ref{tab:tau-diff}, contrast~iv). On the inbound counterparty ratio both coefficients are negative, $-0.128$ ($[-0.135,-0.121]$) for EndorseRank and $-0.019$ ($[-0.029,-0.010]$) for AWP. The ratio divides the distinct other senders by the transfers received from them, so it falls when the same senders return; EndorseRank's negative coefficient says that contracts that many owners approve are, on this measure, contracts whose senders return. The family describes how much longitudinal transfer information each score carries.

\dissertationSubheading[sec:er-awp-holdout]{Out of window}

Table~\ref{tab:endorserank-awp-holdout} collects the paired differences between the two scores in both holdout windows. In every row $a$ is EndorseRank and $b$ AWP.

% Generated by experiments/scripts/export_latex.py -- do not edit by hand.
% Source: experiments/data/2-processed-tables-and-evaluations/revised-usd/holdout-w0/eval_summary.json
% Source: experiments/data/2-processed-tables-and-evaluations/revised-usd/supplementary/supplementary.json
% Source: experiments/data/2-processed-tables-and-evaluations/revised-usd/window-w1/eval_summary.json
\begin{table}[htbp]
\centering
\caption[Paired bootstrap differences between EndorseRank and AWP out of window.]{Paired bootstrap differences between EndorseRank and AWP out of window (W0 spenders $n=17{,}050$; W0 traders $n=88$; W1 spenders $n=18{,}152$; W1 traders $n=70$). None of them enters rule A or rule B: the spender rows were computed after the labels were known, and the trader rows are the EndorseRank-minus-AWP pair of the comparison on an outcome built from neither edge type (Tables~\ref{tab:neutral-label-contrasts} and~\ref{tab:neutral-label-grid}). The W1 spender rows were computed from the saved W1 scores after the W1 evaluation, with the resamples of Table~\ref{tab:fresh-holdout}. Within each block $a$ and $b$ share the contract resamples (400 for the spenders and 2{,}000 for the traders; seed 42).}
\label{tab:endorserank-awp-holdout}
\fitwidth{%
\begin{tabular}{llccccc}
\toprule
Contrast ($a$ minus $b$) & Label & $\tau_a$ & $\tau_b$ & $\Delta\tau$ & 95\% CI of $\Delta\tau$ & $P(\Delta\tau>0)$ \\
\midrule
\multicolumn{7}{l}{\emph{W0, spenders ($n=17{,}050$)}} \\
EndorseRank minus AWP & new approval pairs & 0.231 & 0.196 & +0.035 & [0.019, 0.051] & 1.000 \\
EndorseRank minus AWP & new transfer senders & 0.200 & 0.346 & -0.146 & [-0.156, -0.136] & 0.000 \\
\midrule
\multicolumn{7}{l}{\emph{W0, traders with at least three closes ($n=88$)}} \\
EndorseRank minus AWP & liquidation-free close rate & 0.224 & 0.240 & -0.017 & [-0.156, 0.128] & 0.394 \\
\midrule
\multicolumn{7}{l}{\emph{W1, spenders ($n=18{,}152$)}} \\
EndorseRank minus AWP & new approval pairs & 0.234 & 0.169 & +0.065 & [0.051, 0.081] & 1.000 \\
EndorseRank minus AWP & new transfer senders & 0.194 & 0.330 & -0.136 & [-0.146, -0.127] & 0.000 \\
\midrule
\multicolumn{7}{l}{\emph{W1, traders with at least three closes ($n=70$)}} \\
EndorseRank minus AWP & liquidation-free close rate & 0.297 & 0.329 & -0.032 & [-0.133, 0.066] & 0.279 \\
\bottomrule
\end{tabular}
}
\end{table}


On the spenders of W0, EndorseRank ranks new approval pairs at $\tau=0.241$ (interval $[0.234,0.248]$) against $0.196$ ($[0.187,0.206]$) for AWP, a paired difference of $+0.045$ ($[0.033,0.055]$). On new transfer senders the order reverses: AWP reaches $0.340$ ($[0.333,0.348]$) against $0.226$ ($[0.219,0.234]$) for EndorseRank, a difference of $-0.114$ ($[-0.120,-0.108]$). Every label falls after all scored events, so on this cohort each score ranks the later arrivals of its own edge type better than the other score does. Coverage accounts for part of the first gap. AWP scores zero for the $3{,}753$ spenders without a transfer edge before the freeze, and $11.9\%$ of them gained a new approval pair, against $10.1\%$ of the rest. On the $29{,}228$ spenders that had received a transfer from another address, the order on new approval pairs reverses: AWP reaches $0.262$ ($[0.253,0.270]$) and EndorseRank $0.249$ ($[0.241,0.256]$), without a paired test (Table~\ref{tab:holdout-receiving}). On new senders AWP stays ahead there, at $0.349$ ($[0.342,0.356]$) against $0.247$ ($[0.240,0.255]$).

In July--September, with the contrasts computed from the saved scores of W1 after the registered evaluation, the pattern on the full spender cohort repeats. EndorseRank ranks new approval pairs at $0.233$ ($[0.226,0.240]$) against $0.177$ ($[0.166,0.187]$) for AWP ($+0.056$, $[0.046,0.068]$), and AWP ranks new senders at $0.322$ ($[0.316,0.329]$) against $0.210$ ($[0.203,0.217]$) ($-0.112$, $[-0.119,-0.106]$). On the liquidation-free close rate of the traders, one of the pairs registered with W1, EndorseRank reaches $0.230$ ($[0.115,0.341]$) in W0 and $0.306$ ($[0.177,0.421]$) in W1, and AWP $0.168$ ($[0.022,0.315]$) and $0.240$ ($[0.094,0.365]$). Neither difference can be told from zero ($+0.062$, $[-0.079,0.217]$, and $+0.066$, $[-0.069,0.227]$); the trader cohorts hold $88$ and $70$ contracts, so the intervals are wide.

On the W0 traders' profit labels (Table~\ref{tab:holdout-traders}), every interval of either score on profitable closes and realized gain includes zero: EndorseRank reaches $0.156$ ($[-0.022,0.302]$) and $0.088$ ($[-0.079,0.242]$), AWP $0.138$ ($[-0.020,0.281]$) and $0.134$ ($[-0.018,0.283]$). On the two shares both scores are negative. AWP's intervals lie below zero on both, at $-0.170$ ($[-0.299,-0.032]$) on the profitable share and $-0.253$ ($[-0.384,-0.112]$) on the non-loss share; EndorseRank's lies below zero on the non-loss share only, at $-0.170$ ($[-0.286,-0.044]$), against $-0.053$ ($[-0.186,0.082]$) on the profitable share. The transfer in-degree is negative on both shares as well, at $-0.198$ ($[-0.332,-0.053]$) and $-0.190$ ($[-0.344,-0.036]$). Neither score ranks these contract traders by trading success.

\dissertationSubheading[sec:er-awp-neutral-grid]{All trader labels in both windows}

Table~\ref{tab:neutral-label-grid} sets the allowance side against the transfer side on all six trader labels in both windows. Its four pairs and six labels were registered with W1, before the logs of July--September were extracted (Section~\ref{sub:neutral-label-protocol}). By then the single-score coefficients of the W0 traders on the profit labels were known (Table~\ref{tab:holdout-traders}), while no score had been computed against their liquidation labels. Only P and Q on the liquidation-free close rate of W1 enter a decision, through rules R1 and R2 (Claim~C6, Section~\ref{sub:neutral-label-results}); the column for EndorseRank minus AWP is reported beside them.

% Generated by experiments/scripts/export_latex.py -- do not edit by hand.
% Source: experiments/data/2-processed-tables-and-evaluations/revised-usd/window-w1/eval_summary.json
\begin{table}[htbp]
\centering
\caption[Allowance-side minus transfer-side scores on every trader label in both windows.]{Allowance-side minus transfer-side scores, $\Delta\tau$ with 95\% paired bootstrap interval, on all six trader labels of the window W1 ($n=70$ GMX~V2 contract traders) and of W0 ($n=88$). P: in-approve degree minus transfer in-degree. Q: EndorseRank with activity restarts minus AWP (same edge weights and restart rule). C-PR: $\lambda=1$ minus $\lambda=0$ (same USD build). ER: EndorseRank minus AWP. Profitable closes and realized gain grow with the number of closes; the other four labels do not. CI = 95\% percentile bootstrap over contracts (2{,}000 paired resamples, seed 42).}
\label{tab:neutral-label-grid}
\fitwidth{%
\begin{tabular}{lcccc}
\toprule
Label & P (counts) & Q (PageRanks) & C-PR layers & ER minus AWP \\
\midrule
\multicolumn{5}{l}{\emph{Window W1, July--September 2026 ($n=70$ traders)}} \\
Liquidation-free close rate & \begin{tabular}[c]{@{}c@{}}0.084\\[-1pt]{\footnotesize [-0.013, 0.176]}\end{tabular} & \begin{tabular}[c]{@{}c@{}}-0.041\\[-1pt]{\footnotesize [-0.135, 0.050]}\end{tabular} & \begin{tabular}[c]{@{}c@{}}-0.049\\[-1pt]{\footnotesize [-0.169, 0.060]}\end{tabular} & \begin{tabular}[c]{@{}c@{}}-0.032\\[-1pt]{\footnotesize [-0.133, 0.066]}\end{tabular} \\
No liquidation & \begin{tabular}[c]{@{}c@{}}0.078\\[-1pt]{\footnotesize [-0.012, 0.166]}\end{tabular} & \begin{tabular}[c]{@{}c@{}}-0.039\\[-1pt]{\footnotesize [-0.138, 0.062]}\end{tabular} & \begin{tabular}[c]{@{}c@{}}-0.051\\[-1pt]{\footnotesize [-0.170, 0.064]}\end{tabular} & \begin{tabular}[c]{@{}c@{}}-0.029\\[-1pt]{\footnotesize [-0.136, 0.080]}\end{tabular} \\
Profitable closes & \begin{tabular}[c]{@{}c@{}}-0.143\\[-1pt]{\footnotesize [-0.318, 0.013]}\end{tabular} & \begin{tabular}[c]{@{}c@{}}-0.231\\[-1pt]{\footnotesize [-0.429, -0.027]}\end{tabular} & \begin{tabular}[c]{@{}c@{}}-0.190\\[-1pt]{\footnotesize [-0.411, 0.037]}\end{tabular} & \begin{tabular}[c]{@{}c@{}}-0.214\\[-1pt]{\footnotesize [-0.423, -0.012]}\end{tabular} \\
Realized gain & \begin{tabular}[c]{@{}c@{}}-0.135\\[-1pt]{\footnotesize [-0.308, 0.025]}\end{tabular} & \begin{tabular}[c]{@{}c@{}}-0.350\\[-1pt]{\footnotesize [-0.536, -0.157]}\end{tabular} & \begin{tabular}[c]{@{}c@{}}-0.404\\[-1pt]{\footnotesize [-0.624, -0.182]}\end{tabular} & \begin{tabular}[c]{@{}c@{}}-0.346\\[-1pt]{\footnotesize [-0.547, -0.144]}\end{tabular} \\
Profitable share & \begin{tabular}[c]{@{}c@{}}-0.196\\[-1pt]{\footnotesize [-0.364, -0.046]}\end{tabular} & \begin{tabular}[c]{@{}c@{}}-0.336\\[-1pt]{\footnotesize [-0.523, -0.141]}\end{tabular} & \begin{tabular}[c]{@{}c@{}}-0.344\\[-1pt]{\footnotesize [-0.551, -0.120]}\end{tabular} & \begin{tabular}[c]{@{}c@{}}-0.316\\[-1pt]{\footnotesize [-0.517, -0.112]}\end{tabular} \\
Non-loss share & \begin{tabular}[c]{@{}c@{}}-0.252\\[-1pt]{\footnotesize [-0.391, -0.120]}\end{tabular} & \begin{tabular}[c]{@{}c@{}}-0.350\\[-1pt]{\footnotesize [-0.526, -0.175]}\end{tabular} & \begin{tabular}[c]{@{}c@{}}-0.324\\[-1pt]{\footnotesize [-0.531, -0.109]}\end{tabular} & \begin{tabular}[c]{@{}c@{}}-0.346\\[-1pt]{\footnotesize [-0.526, -0.161]}\end{tabular} \\
\midrule
\multicolumn{5}{l}{\emph{Window W0, April--June 2026 ($n=88$ traders)}} \\
Liquidation-free close rate & \begin{tabular}[c]{@{}c@{}}0.182\\[-1pt]{\footnotesize [0.052, 0.314]}\end{tabular} & \begin{tabular}[c]{@{}c@{}}-0.014\\[-1pt]{\footnotesize [-0.147, 0.130]}\end{tabular} & \begin{tabular}[c]{@{}c@{}}-0.041\\[-1pt]{\footnotesize [-0.187, 0.105]}\end{tabular} & \begin{tabular}[c]{@{}c@{}}-0.017\\[-1pt]{\footnotesize [-0.156, 0.128]}\end{tabular} \\
No liquidation & \begin{tabular}[c]{@{}c@{}}0.197\\[-1pt]{\footnotesize [0.059, 0.333]}\end{tabular} & \begin{tabular}[c]{@{}c@{}}-0.010\\[-1pt]{\footnotesize [-0.152, 0.146]}\end{tabular} & \begin{tabular}[c]{@{}c@{}}-0.051\\[-1pt]{\footnotesize [-0.198, 0.102]}\end{tabular} & \begin{tabular}[c]{@{}c@{}}-0.012\\[-1pt]{\footnotesize [-0.154, 0.137]}\end{tabular} \\
Profitable closes & \begin{tabular}[c]{@{}c@{}}0.075\\[-1pt]{\footnotesize [-0.080, 0.222]}\end{tabular} & \begin{tabular}[c]{@{}c@{}}-0.016\\[-1pt]{\footnotesize [-0.150, 0.118]}\end{tabular} & \begin{tabular}[c]{@{}c@{}}-0.067\\[-1pt]{\footnotesize [-0.222, 0.087]}\end{tabular} & \begin{tabular}[c]{@{}c@{}}-0.005\\[-1pt]{\footnotesize [-0.163, 0.155]}\end{tabular} \\
Realized gain & \begin{tabular}[c]{@{}c@{}}0.017\\[-1pt]{\footnotesize [-0.150, 0.176]}\end{tabular} & \begin{tabular}[c]{@{}c@{}}-0.107\\[-1pt]{\footnotesize [-0.252, 0.044]}\end{tabular} & \begin{tabular}[c]{@{}c@{}}-0.248\\[-1pt]{\footnotesize [-0.433, -0.067]}\end{tabular} & \begin{tabular}[c]{@{}c@{}}-0.118\\[-1pt]{\footnotesize [-0.286, 0.046]}\end{tabular} \\
Profitable share & \begin{tabular}[c]{@{}c@{}}0.045\\[-1pt]{\footnotesize [-0.113, 0.184]}\end{tabular} & \begin{tabular}[c]{@{}c@{}}0.046\\[-1pt]{\footnotesize [-0.071, 0.172]}\end{tabular} & \begin{tabular}[c]{@{}c@{}}0.027\\[-1pt]{\footnotesize [-0.112, 0.181]}\end{tabular} & \begin{tabular}[c]{@{}c@{}}0.069\\[-1pt]{\footnotesize [-0.060, 0.208]}\end{tabular} \\
Non-loss share & \begin{tabular}[c]{@{}c@{}}-0.080\\[-1pt]{\footnotesize [-0.228, 0.058]}\end{tabular} & \begin{tabular}[c]{@{}c@{}}0.055\\[-1pt]{\footnotesize [-0.070, 0.192]}\end{tabular} & \begin{tabular}[c]{@{}c@{}}0.037\\[-1pt]{\footnotesize [-0.111, 0.189]}\end{tabular} & \begin{tabular}[c]{@{}c@{}}0.076\\[-1pt]{\footnotesize [-0.063, 0.223]}\end{tabular} \\
\bottomrule
\end{tabular}
}
\end{table}


In W1 no pair separates the two sides on either liquidation label: every interval includes zero, P on the liquidation-free close rate at $+0.063$ ($[-0.026,0.148]$) and the EndorseRank pair at $+0.066$ ($[-0.069,0.227]$). In W0 the counts favor the allowance side on both liquidation labels, P at $+0.252$ ($[0.093,0.413]$) on the liquidation-free close rate and $+0.270$ ($[0.098,0.441]$) on the no-liquidation flag, while the three PageRank pairs lean the same way with intervals that include zero. On the four profit labels of W1 every pair is negative. The intervals of Q and of the EndorseRank pair lie below zero on all four, those of P on the two shares, and those of the C-PR layers on realized gain and on both shares. In W0 most intervals on the profit labels include zero, and Q and the EndorseRank pair lean to the allowance side on the profitable share ($+0.093$, $[0.001,0.187]$, and $+0.117$, $[0.004,0.227]$), so the profit labels show no direction that holds in both windows. With $70$ and $88$ traders, none of these differences is precise.
